\section{Introduction}
\label{sec:intro}


% High-fidelity 3D reconstruction from monocular image sequences remains a critical challenge in computer vision, with broad applications in AR/VR, robotics, and real-to-sim environments. Recently, the integration of 3D Gaussian Splatting (3DGS)~\cite{kerbl20233d} with deep visual SLAM has driven significant progress, enabling real-time rendering and direct geometric manipulation. To improve robustness in unconstrained streams, recent 3DGS-based SLAM frameworks~\cite{liartdeco, cheng2025outdoor} have increasingly incorporated foundation models as geometric priors.

% 3D scene reconstruction has become a fundamental capability in computer vision, enabling a wide range of applications including embodied intelligence, robotic perception, autonomous navigation, and large-scale scene digitization~\cite{zhou2025hugsim,huang2026soma,yu2026gaussexplorer,yang2025novel}. 
% In recent years, the field has advanced rapidly with the emergence of neural scene representations and visual foundation models. 
% Per-scene optimization methods such as 3D Gaussian Splatting (3DGS)~\cite{kerbl20233d} achieve high-fidelity reconstruction with real-time rendering~\cite{lu2024scaffold,yu2024mip,ren2024octree}, while feed-forward 3D reconstruction models~\cite{leroy2024grounding,wang2025pi,wang2025vggt} infer dense geometric priors such as camera poses and depth maps from images in a single inference pass. 
% Despite their strong performance, most existing methods operate in an \emph{offline} paradigm, where all images are processed jointly after acquisition. 
% This constraint inhibits real-time feedback and limits scalability in long video streams or continuously expanding scenes.
% Consequently, there is a pressing need for streaming reconstruction, where camera poses and scene representations are continuously updated as new frames arrive, enabling immediate feedback and aligning naturally with SLAM frameworks for real-time perception.

% Existing efforts toward \emph{streaming} 3D reconstruction can be broadly categorized into two families. The first family extends feed-forward reconstruction models with a memory mechanism that summarizes previously processed frames and interacts with incoming observations to predict geometry incrementally~\cite{spann3r,long3r,point3r}. 
% While these methods are efficient and naturally streaming, they typically operate on reduced-resolution and produce inaccurate 3D reconstruction. The second family integrates foundation-model priors into a SLAM pipeline to guide pose estimation and reconstruction~\cite{murai2024_mast3rslam,artdeco,maggio2025vggt,maggio2026vggt}. However, existing approaches face a fundamental trade-off: \emph{either} they rely on pairwise priors (e.g., MASt3R-SLAM~\cite{murai2024_mast3rslam}), requiring multiple model inferences for each incoming frame and introducing redundant computation; \emph{or} they adopt multi-frame priors (e.g., VGGT-SLAM 2.0~\cite{maggio2026vggt}), which lack pixel-level dense correspondences and therefore cannot fully exploit the geometric constraints of the SLAM framework.

% % limitation 分析中，matching的重要性更早一点提出， 篇幅加大一点。比如指出xxx情况下普遍会fail

% % 然后tune foundation model这件事写成一个不得不解决的需求，因为当前最强的model的关注点在出几何结构， 然后视角间关联能力是普遍缺乏的。所以急需通过后训练来强化这部分功能。

% To address these limitations, we propose \textbf{M$^3$}, a streaming 3D reconstruction framework that tightly integrates a multi-view feed-forward foundation model with a SLAM pipeline. We first enhance a state-of-the-art multi-view geometric foundation model by introducing a dense matching head that enables pixel-level correspondences across frames, allowing the SLAM framework to fully exploit geometric constraints for accurate pose refinement and scene reconstruction. The enhanced model is then tightly incorporated into the SLAM framework, where each update is performed with a single feed-forward inference over both historical keyframes and incoming frames. This design not only improves reconstruction and pose estimation quality but also significantly reduces redundant model invocations, leading to a more efficient streaming pipeline. In addition, we introduce a dynamic region identification module to detect and suppress transient objects, improving the robustness of the framework in real-world scenarios.
% \jcj{Extensive experiments across diverse indoor and outdoor benchmarks show that M$^3$ achieves state-of-the-art accuracy in both pose estimation and 3D reconstruction, while maintaining competitive efficiency on long monocular video streams.}

3D scene reconstruction has become a fundamental capability in computer vision, enabling applications ranging from robotic perception to large-scale scene digitization~\cite{zhou2025hugsim,huang2026soma,yu2026gaussexplorer,yang2025novel}. Recently, the field has been revolutionized by two paradigms: per-scene optimization, such as 3D Gaussian Splatting (3DGS)~\cite{kerbl20233d}, which delivers high-fidelity rendering, and feed-forward geometric foundation models~\cite{leroy2024grounding,wang2025pi,wang2025vggt}, which infer dense priors in a single pass. However, most existing foundation models are inherently batch-oriented, designed to process a fixed set of images jointly. This offline nature precludes real-time feedback and limits scalability in open-ended environments, underscoring the urgent need for streaming reconstruction, where camera trajectories and scene geometry are incrementally updated as new observations arrive.

Existing efforts toward streaming 3D reconstruction generally follow two trajectories, yet both face significant hurdles. The first family attempts to adapt feed-forward models to a streaming context by incorporating memory mechanisms that summarize past observations to predict geometry incrementally~\cite{spann3r,long3r,point3r}. While these methods are efficient, they typically produce low-resolution results and struggle with cumulative drift, as they lack the iterative global refinement mechanisms found in classical SLAM. The second family instead integrates foundation-model priors into a SLAM pipeline to guide optimization~\cite{murai2024_mast3rslam,artdeco,maggio2025vggt}. However, these approaches are often trapped in a fundamental trade-off: pairwise-prior methods, such as MASt3R-SLAM~\cite{murai2024_mast3rslam}, suffer from redundant computation and quadratic complexity, whereas multi-frame prior methods like VGGT-SLAM 2.0~\cite{maggio2026vggt} provide global geometry but lack the pixel-level dense correspondences necessary for rigorous geometric optimization.

We argue that the primary bottleneck in current multi-view foundation models is their disproportionate focus on individual scene geometry at the expense of inter-view relational consistency. While these models produce impressive 3D structures, they are often blind to the precise pixel-to-pixel associations across frames. Without such fine-grained correspondences, the SLAM backend cannot establish the strong epipolar constraints required for Bundle Adjustment (BA), leading to catastrophic failures like ghosting artifacts or trajectory divergence in complex sequences. Consequently, fine-tuning foundation models specifically to recover dense matching is no longer an option, but a necessity to unlock their full potential for downstream SLAM tasks.

To bridge this gap, we propose M$^3$, a streaming 3D reconstruction framework that tightly couples a multi-view foundation model with a robust SLAM pipeline. Our approach first enhances a state-of-the-art multi-view geometric foundation model by introducing a dedicated dense matching head, specifically trained to recover pixel-level correspondences. This enables the SLAM framework to leverage the foundation model’s geometry for accurate, high-frequency pose refinement. Unlike previous black-box integrations, M$^3$ performs a single feed-forward inference over both historical keyframes and incoming frames to simultaneously update geometry and tracking, significantly reducing redundant model invocations. Furthermore, we introduce a dynamic region identification module to detect and suppress transient objects, ensuring stable static scene reconstruction in real-world environments. Extensive experiments across diverse indoor and outdoor benchmarks demonstrate that M$^3$ achieves state-of-the-art accuracy in both pose estimation and 3D reconstruction, while maintaining competitive efficiency on long-duration monocular video streams.

In summary, our core contributions are as follows:

\begin{itemize}
    \item We introduce a dedicated matching head to a multi-view foundation model, leveraging pixel-level descriptors to facilitate refined cross-frame dense matching for rigorous geometric optimization.
    \item We propose M$^3$, a SLAM framework leveraging a multi-view foundation model to simultaneously facilitate frontend tracking and backend global optimization via a single feed-forward inference.
    \item Extensive experiments across diverse benchmarks demonstrate that M$^3$ delivers state-of-the-art accuracy in both pose estimation and 3D reconstruction, while maintaining high computational efficiency on long-duration monocular sequences.
\end{itemize}

% However, current 3D foundation model-based SLAM frameworks fail to fully exploit the inherent multi-view consistency offered by these large-scale priors. Existing frameworks typically fall into one of two sub-optimal paradigms. The first paradigm relies on models restricted to pairwise image inference, which inherently cannot capture global multi-view consistency across a broader sequence in a single pass. This forces the framework to artificially decouple frontend tracking from backend optimization, necessitating complex, redundant fusion modules to stitch together multiple pairwise inferences into a globally consistent map. The second paradigm employs multi-frame inference but lacks explicit, dense pixel-level correspondence modeling. Without accurate explicit correspondences, it becomes nearly impossible to seamlessly integrate classical visual SLAM optimization techniques, thereby severely bottlenecking the geometric precision of the reconstruction. Furthermore, unconstrained real-world videos are inevitably populated with dynamic objects and transient elements. Most existing pipelines blindly integrate these regions into the static scene representation, introducing severe geometric artifacts, corrupting camera pose estimates, and leading to catastrophic trajectory drift.

% To address these critical bottlenecks, we propose M$^3$, an efficient and robust streaming reconstruction framework designed for uncalibrated, arbitrary monocular videos. Specifically, to better leveraging multi-view consistency, we augment a multi-frame 3D foundation model (Pi3X) with a dense matching head. This advancement enables the network to simultaneously predict multi-frame geometric priors and dense, pixel-level feature descriptors in a single forward pass. Leveraging these explicit dense correspondences, we employ a sliding window mechanism to successfully unify frontend tracking and backend global optimization within a tightly coupled architecture. By jointly processing historical and incoming frames in a single multi-view inference, our approach not only eliminates the substantial computational redundancy of frequent pairwise inferences but also seamlessly integrates classical visual SLAM optimization techniques, fundamentally ensuring strict multi-frame global consistency. Furthermore, to tackle the issue of dynamic objects corrupting pose tracking and polluting the reconstructed map, we design a descriptor-based motion estimation module. By computing the descriptor consistency between historical keyframes warped into the current view, this module explicitly generates a precise motion map to localize transient elements. Utilizing this mask, the framework actively filters or significantly down-weights dynamic regions during joint camera pose optimization on the $Sim(3)$ manifold and 3D Gaussian primitive initialization. This effectively isolates interference from complex real-world environments, radically suppressing trajectory drift and rendering artifacts.

% Our core improvements and contributions are summarized as follows:


% \begin{itemize}
%     \item \textbf{Unified Framework with Augmented Multi-View Priors}: We augment a multi-frame foundation model (Pi3X) with a dense matching head, effectively bridging the gap between multi-view inference and explicit pixel-level correspondences. This advancement allows us to tightly couple frontend tracking and backend global bundle adjustment via a sliding window mechanism, eliminating redundant pairwise inferences while natively enforcing multi-view consistency.
%     \item \textbf{Explicit Dynamic Region Filtering}: To tackle the challenges of complex, unconstrained real-world environments, we introduce a descriptor-based motion estimation module. By computing descriptor consistency across warped frames, the framework predicts a precise motion map to explicitly identify and down-weight dynamic regions, significantly reducing trajectory drift and geometric artifacts during 3D Gaussian initialization.
%     \item \textbf{Robust $Sim(3)$ Optimization and Hierarchical Loop Closure}: To prevent the accumulation of scale drift inherent in sequential uncalibrated foundation model inferences, we formulate camera pose tracking and global optimization within the $Sim(3)$ manifold. Coupled with a hierarchical loop closure strategy—utilizing SALAD for rapid candidate retrieval and the augmented Pi3X for dense geometric validation—our framework ensures highly accurate tracking and globally consistent reconstructions.
% \end{itemize}

% The typical SLAM framework consists of two fundamental components: a tracking module (Front-end) and a global optimization module (Back-end). The former focuses on low-latency tracking of the current camera pose and local geometry, while the latter enforces global consistency across the entire trajectory and reconstructed map. Prior works, such as MASt3R-SLAM~\cite{}, typically decouple these modules, performing independent inference passes for local tracking and global refinement. This separation is largely necessitated by the limitations of the underlying foundation models, which are restricted to pairwise image processing and cannot handle multi-view batches efficiently. 
